% !TeX program = xelatex
\documentclass[11pt,a4paper]{ctexart}

\usepackage[left=2.4cm,right=2.4cm,top=2.1cm,bottom=2.1cm]{geometry}
\usepackage{amsmath,amssymb,mathtools,braket}
\usepackage{xcolor}
\usepackage[inline]{enumitem}
\usepackage{fancyhdr}
\usepackage[hidelinks]{hyperref}

% ============================================================
% 基本信息
% ============================================================
\newcommand{\coursename}{高等量子力学}
\newcommand{\assignmentnumber}{第 4 次作业}
\newcommand{\teachername}{\underline{\makebox[2.8cm][c]{金彪}}}
\newcommand{\studentname}{\underline{\makebox[2.8cm][c]{林照翔}}}
\newcommand{\studentid}{\underline{\makebox[3.0cm][c]{202628001207046}}}

% ============================================================
% 页面样式
% ============================================================
\setlength{\parindent}{2em}
\setlength{\parskip}{0.25em}
\setlength{\headheight}{14pt}
\pagestyle{fancy}
\fancyhf{}
\lhead{\coursename}
\rhead{\assignmentnumber}
\cfoot{\thepage}
\renewcommand{\headrulewidth}{0.4pt}

% ============================================================
% 列表参数
% ============================================================
\setlist[enumerate,1]{
  label=(\alph*),
  leftmargin=3.2em,
  itemsep=0.55em,
  topsep=0.45em
}

\newlist{inlineenum}{enumerate*}{1}
\setlist[inlineenum,1]{
  label=(\alph*),
  labelsep=0.5em,
  itemjoin={；\quad},
  itemjoin*={；\quad}
}

% ============================================================
% 题目与回答空间
% ============================================================
\newcounter{problem}
\newenvironment{problem}
  {\refstepcounter{problem}%
   \par\medskip\noindent
   \begingroup
   \setlength{\fboxsep}{3pt}%
   \setlength{\fboxrule}{0.8pt}%
   \fbox{\bfseries 题目\chinese{problem}}%
   \endgroup
   \quad\ignorespaces}
  {\par\medskip}

\newenvironment{answer}
  {\par\smallskip\noindent
   \begingroup
   \setlength{\fboxsep}{4pt}%
   \colorbox{black}{\textcolor{white}{\bfseries 回答}}%
   \endgroup
   \quad\ignorespaces}
  {\par\smallskip\noindent\rule{\linewidth}{0.4pt}\medskip}

\newcommand{\answerspace}[1]{\par\vspace*{#1}}

\begin{document}

\begin{center}
  {\LARGE\bfseries \coursename\quad \assignmentnumber}\par
  \vspace{0.7em}
  \small
  教师：\teachername\qquad
  姓名：\studentname\qquad
  学号：\studentid
\end{center}
\vspace{0.4em}

% 本份作业承接原题前二十题，题号从二十一开始。
\setcounter{problem}{20}

\begin{problem}
课堂笔记内容：Schmidt 分解部分。
\begin{equation}
  \hat\rho_{\mathrm{R}}(1)
  =\sum_{i=1}^{N}\rho_i\Ket{\tilde u_i(1)}\Bra{\tilde u_i(1)},
  \qquad
  \hat\rho_{\mathrm{R}}(2)
  =\sum_{i=1}^{N}\rho_i\Ket{\tilde v_i(2)}\Bra{\tilde v_i(2)}.
\end{equation}
证明当 $N\geq 2$ 时，约化密度算符 $\hat\rho_{\mathrm{R}}(1)$ 及
$\hat\rho_{\mathrm{R}}(2)$ 皆为混合态密度算符。
\end{problem}

\begin{answer}
\answerspace{13cm}
\end{answer}
\clearpage

\begin{problem}
\begin{enumerate}[label=(\arabic*)]
  \item 试判断以下密度矩阵描述的是纯态还是混合态。
  \begin{equation}
    \rho_1=\frac12\begin{pmatrix}1&0\\0&1\end{pmatrix};\quad
    \rho_2=\begin{pmatrix}1&0\\0&0\end{pmatrix};\quad
    \rho_3=\frac12\begin{pmatrix}1&1\\1&1\end{pmatrix};\quad
    \rho_4=\frac13\begin{pmatrix}2&1\\1&1\end{pmatrix}.
  \end{equation}
  \item 求算符
  $A=\begin{pmatrix}1&1-\mathrm{i}\\1+\mathrm{i}&-1\end{pmatrix}$
  在这些态上的平均值。
  \item 求这些密度矩阵的 von Neumann 熵。
\end{enumerate}
\end{problem}

\begin{answer}
\answerspace{15cm}
\end{answer}
\clearpage

\begin{problem}
设两电子自旋系统状态为
\begin{equation}
  \begin{aligned}
    \Ket{\psi(1,2)}
    ={}&\frac{1}{\sqrt2}\Ket{\uparrow}_1\otimes
    \left(\frac12\Ket{\uparrow}_2+\frac{\sqrt3}{2}\Ket{\downarrow}_2\right)\\
    &+\frac{1}{\sqrt2}\Ket{\downarrow}_1\otimes
    \left(\frac{\sqrt3}{2}\Ket{\uparrow}_2+\frac12\Ket{\downarrow}_2\right).
  \end{aligned}
\end{equation}
\begin{enumerate}[label=(\arabic*)]
  \item 进行 Schmidt 分解。
  \item 求约化密度矩阵 $\rho_1,\rho_2$。
\end{enumerate}
\end{problem}

\begin{answer}
\answerspace{15cm}
\end{answer}
\clearpage

\begin{problem}
取两电子自旋系统基矢为
\begin{equation}
  \begin{aligned}
    \Ket{\chi_1}&=\Ket{\uparrow}_1\otimes\Ket{\uparrow}_2,&
    \Ket{\chi_2}&=\Ket{\uparrow}_1\otimes\Ket{\downarrow}_2,\\
    \Ket{\chi_3}&=\Ket{\downarrow}_1\otimes\Ket{\uparrow}_2,&
    \Ket{\chi_4}&=\Ket{\downarrow}_1\otimes\Ket{\downarrow}_2.
  \end{aligned}
\end{equation}
设纯态密度矩阵为
\begin{equation}
  \rho(1,2)=
  \begin{pmatrix}
    \dfrac12&\dfrac{\sqrt2}{4}&0&\dfrac{\sqrt2}{4}\\[0.5em]
    \dfrac{\sqrt2}{4}&\dfrac14&0&\dfrac14\\[0.5em]
    0&0&0&0\\[0.5em]
    \dfrac{\sqrt2}{4}&\dfrac14&0&\dfrac14
  \end{pmatrix}.
\end{equation}
求约化密度矩阵 $\rho_1,\rho_2$。
\end{problem}

\begin{answer}
\answerspace{13cm}
\end{answer}
\clearpage

\begin{problem}
课堂例题（绝热近似部分）。求：
\begin{enumerate}
  \item $\nabla_B H\bigl(\vec B(t)\bigr)$；
  \item $\hat A_+\bigl(\vec B(t)\bigr)
    =\mathrm{i}\Braket{\chi_+,\vec B(t)|\nabla_B|\chi_+,\vec B(t)}$；
  \item $\hat B_+\bigl(\vec B(t)\bigr)
    =\nabla_B\times\hat A_+\bigl(\vec B(t)\bigr)$；
  \item Berry 相位
  \begin{equation}
    \gamma_+(T)=\iint_S\hat B_+\bigl(\vec B(T)\bigr)\cdot\mathrm{d}\vec S,
    \qquad T=\frac{2\pi}{\omega}.
  \end{equation}
\end{enumerate}
\end{problem}

\begin{answer}
\answerspace{15cm}
\end{answer}

\end{document}
